\subsection{The max-type threshold test}
\label{sec:results-stageB}

In the null calibration both Stage~B rules keep the false-split rate below
0.05 in all 36 design cells (18 task, sample size, and bin settings under two
stopping modes), but that design cannot expose a failure of the
threshold test. Every split must first pass the Stage~A test over five
predictors, so even without the threshold correction the rate stays between
0.0355 and 0.049 (\cref{app:stageB-tables}). The Stage-B-only study therefore
isolates the threshold test on one predictor fixed without the response, with
the Stage~A test disabled (\cref{tab:stageB-only}).

\begin{table}[tb]
  \centering
  \footnotesize
  \setlength{\tabcolsep}{4pt}
  \begin{tabular}{@{}lrrrrrrrr@{}}
    \toprule
     & & & \multicolumn{3}{c}{Size at nominal 0.05} & \multicolumn{3}{c}{Size-matched power} \\
    \cmidrule(lr){4-6}\cmidrule(l){7-9}
    \tablehead{Task} & $n$ & Bins & Per-threshold & Max-type & No correction & Match size & Per-threshold & Max-type \\
    \midrule
    Classification & 100 & 16 & 0.013 & 0.032 & 0.156 & 0.050 & 0.390 & 0.379 \\
     & 100 & 64 & 0.007 & 0.036 & 0.218 & 0.044 & 0.385 & 0.354 \\
     & 100 & 256 & 0.007 & 0.036 & 0.226 & 0.034 & 0.335 & 0.318 \\
     & 200 & 16 & 0.010 & 0.037 & 0.152 & 0.050 & 0.643 & 0.618 \\
     & 200 & 64 & 0.008 & 0.035 & 0.242 & 0.050 & 0.630 & 0.615 \\
     & 200 & 256 & 0.003 & 0.039 & 0.298 & 0.031 & 0.562 & 0.538 \\
     & 500 & 16 & 0.015 & 0.032 & 0.172 & 0.050 & 0.958 & 0.953 \\
     & 500 & 64 & 0.007 & 0.029 & 0.270 & 0.050 & 0.961 & 0.964 \\
     & 500 & 256 & 0.000 & 0.023 & 0.314 & 0.034 & 0.958 & 0.952 \\
    Regression & 100 & 16 & 0.019 & 0.023 & 0.210 & 0.050 & 0.292 & 0.281 \\
     & 100 & 64 & 0.013 & 0.029 & 0.282 & 0.050 & 0.305 & 0.285 \\
     & 100 & 256 & 0.011 & 0.031 & 0.286 & 0.050 & 0.317 & 0.299 \\
     & 200 & 16 & 0.027 & 0.033 & 0.230 & 0.050 & 0.591 & 0.581 \\
     & 200 & 64 & 0.017 & 0.033 & 0.298 & 0.050 & 0.619 & 0.598 \\
     & 200 & 256 & 0.008 & 0.041 & 0.324 & 0.050 & 0.632 & 0.579 \\
     & 500 & 16 & 0.026 & 0.031 & 0.236 & 0.050 & 0.947 & 0.942 \\
     & 500 & 64 & 0.015 & 0.029 & 0.334 & 0.050 & 0.966 & 0.952 \\
     & 500 & 256 & 0.005 & 0.026 & 0.384 & 0.046 & 0.963 & 0.951 \\
    \bottomrule
  \end{tabular}
  \caption{Stage-B-only study: the threshold test in isolation, with one
  Gaussian predictor fixed without the response and the Stage~A test disabled.
  Size is the false-split rate of a depth-one tree at nominal level 0.05 under
  adaptive stopping, over 1,500 null replicates per cell (three seeds of 500)
  for the two rules, with every 95\% Clopper-Pearson upper limit below 0.056,
  and over 500 for No correction, the per-threshold rule with each candidate
  tested at 0.05. Sizes under no stopping agree to within 0.010 (exactly for
  No correction). Size-matched power under adaptive stopping at the middle effect ($\delta = 0.10$ in
  classification, $0.4$ in regression) is interpolated against realized size
  over the nominal grid from 0.02 to 0.50 (600 replicates per cell and level
  from 0.02 to 0.20, and 200 at $n=100$ and 400 at $n=200$ and 500 at 0.35 and
  0.50) to the match size: 0.05, or the largest realized size both rules reach
  where the per-threshold rule does not reach 0.05.}
  \label{tab:stageB-only}
\end{table}

Both rules are valid at a fixed node, and they differ in how their size
depends on the candidate count. The per-threshold size falls as $K$ grows, with
means 0.018, 0.011, and 0.006 at 16, 64, and 256 bins over both stopping modes,
which is the reverse cardinality bias of \cref{sec:rank-construction} at a
single node; the max-type size stays near 0.035 (means 0.033, 0.035, and
0.035). Without the threshold correction the size is 0.152 to 0.384, 3 to 8
times the nominal level, and it grows with the candidate count at every sample
size in both tasks.

At matched realized size the max-type rule has slightly lower power. The
max-type minus per-threshold difference averages $-0.006$ (95\% interval
$-0.011$ to $-0.002$ from resampling the 18 cells, which treats each matched
power as exact; \cref{app:stageB-tables} gives the simulation error of each
difference); the max-type rule is lower
in 67 and higher in 25 of the 108 cell and effect combinations, and the mean is
$-0.005$ ($-0.009$ to $-0.002$) over the 13 cells matched at size 0.05. Two reruns give $-0.006$ (a fourth
seed, on the 16 cells it covers) and $-0.004$ (the step at the 85th percentile). At a
common nominal level of 0.05 the max-type rule splits more often, by 0.03 on
average and up to 0.12 in the power study (0.015, 0.029, and 0.049 at 16, 64,
and 256 bins) and by 0.07 on average in the Stage-B-only study, a gap set by
the lower realized size of the per-threshold rule.

The max-type rule is cheaper once a node has many candidates. The
head-to-head forests, fitted at the minimum budget, show the largest savings:
the max-type rule is 10 to 27 times faster at $n\le 8{,}000$ and 21 to 36 times
faster at $n=32{,}000$ on the synthetic cells (13 to 36 times on the cells that
\citet{MilletichDownesHirst2026citrees} tabulate), against 1.6 to 7 times on
the runtime panel (\cref{tab:runtime-ablations}), whose nodes admit fewer
candidates. For one tree with no stopping the ratio grows with the bin setting,
as the $K^2/\alpha$ against $BK$ accounting of \cref{sec:stageA-stageB}
predicts: the max-type rule is 1.1 to 1.3 times slower at 16 bins and 1.5 to
2.1 and 2.8 to 12 times faster at 64 and 256 bins (\cref{tab:stageB-scaling}).

On the six Stage~B real datasets the max-type rule scores higher in 8 of the
12 dataset and stopping-mode cells, equal in one, and lower in the other three,
and it fits faster on the four datasets whose nodes admit many candidates
(\cref{app:stageB-tables}). The Spearman
agreement of the two rules' feature importances is 0.36 to 0.88, so the
max-type rule produces a different ranking rather than the same ranking at
lower cost. On the benchmark that different ranking is slightly worse. Standing in for its
per-threshold counterpart on the all-method panels, the max-type forest ranks
lower in both tasks and the max-type tree keeps its position (max-type rows of
\cref{tab:method-main-comparison}). The paired max-type minus per-threshold
score, within each dataset over downstream learners and standard $k$, is
$-0.0014$ and $-0.0001$ for CIF and $-0.0022$ and $-0.0085$ for CIT in
classification and regression (over 22, 8, 23, and 8 datasets), and the
per-threshold rule is the higher of the two on 59 to 63\% of datasets.
\Cref{app:methods} accounts for the ranking runs, seven of which were lost to a
host fault.
